% from ega0-22-fr.tex
\oldpage[0\textsubscript{IV-1}]{210}

une base de
$(\mathfrak qA_{\mathfrak p}/\mathfrak q^2A_{\mathfrak p})
\otimes_{A_{\mathfrak p}}L$~; alors les $F_i$ forment une base de l'image de
$j$~; mais on sait (21.9.3) que $\widehat\Omega^1_{A/k}$ est un $A$-module
libre de rang $r$ et les $k$-dérivations de $A$ dans lui-même engendrent son
dual~; le fait que les $F_i$ soient linéairement indépendants entraîne donc
l'existence de $k$-dérivations $D_j$ de $A$ dans lui-même telles que la matrice
$(\langle F_i,D_j\otimes1\rangle)$ soit inversible, autrement dit la condition
b).

Ces remarques montrent déjà que b) entraîne a). Pour prouver inversement que
la condition a) entraîne que $j$ est injectif, considérons le diagramme
commutatif
\begin{equation}
\label{0.22.7.2.3-fr}
\xymatrix{
  (\mathfrak q/\mathfrak q^2)\otimes_AL
    \ar[r]^j \ar[d]_\alpha &
  \widehat\Omega^1_{A/k}\otimes_AL \ar@{=}[d] \\
  (\mathfrak p/\mathfrak p^2)\otimes_AL
    \ar[r]_i &
  \widehat\Omega^1_{A/k}\otimes_AL
}
\tag{22.7.2.3}
\end{equation}
et notons le lemme suivant~:
\end{proof}

\begin{lemma}[22.7.2.4]
\label{0.22.7.2.4-fr}
Soient $R$ un anneau local régulier, $\mathfrak m$ son idéal maximal,
$K=R/\mathfrak m$ son corps résiduel. Si $\mathfrak n$ est un idéal de $R$ tel
que $R/\mathfrak n$ soit régulier, alors l'homomorphisme canonique
\begin{equation}
\label{0.22.7.2.5-fr}
  \alpha:(\mathfrak n/\mathfrak n^2)\otimes_RK
  \longrightarrow(\mathfrak m/\mathfrak m^2)\otimes_RK
  \tag{22.7.2.5}
\end{equation}
est injectif.
\end{lemma}

\begin{proof}
En effet, on sait (17.1.6) que $\mathfrak n$ est engendré par une suite
$(x_i)_{1\leq i\leq t}$ faisant partie d'un système régulier de paramètres
$(x_i)_{1\leq i\leq n}$ de $R$~; les images des $x_i$ $(1\leq i\leq t)$
forment une base de $(\mathfrak n/\mathfrak n^2)\otimes_RK$, et leurs images
par $\alpha$ font partie de la base de
$(\mathfrak m/\mathfrak m^2)\otimes_RK$ formée des images des $x_i$ pour
$1\leq i\leq n$, d'où le lemme.
\end{proof}

\begin{proof}[Suite de la démonstration de la proposition~{\rm(22.7.2)}]
Ce lemme et le diagramme (22.7.2.3) ramènent donc (en vertu de l'hypothèse a))
à prouver que $i$ est injectif. Or, dans la suite
\begin{equation}
\label{0.22.7.2.6-fr}
  \mathfrak p/\mathfrak p^2\xrightarrow{h}
  \widehat\Omega^1_{A/k}\widehat\otimes_A(A/\mathfrak p)
  \xrightarrow{g}\widehat\Omega^1_{(A/\mathfrak p)/k}\longrightarrow0
  \tag{22.7.2.6}
\end{equation}
on sait (20.7.20) que $g$ est surjectif et que l'image de $h$ est dense dans
le noyau de $g$ pour la topologie $\mathfrak m$-adique ($\mathfrak m$ étant
l'idéal maximal de $A$)~; comme $\widehat\Omega^1_{A/k}$ est un $A$-module de
type fini, tous les sous-modules du $(A/\mathfrak p)$-module
\[
  \widehat\Omega^1_{A/k}\otimes_A(A/\mathfrak p)
  =\widehat\Omega^1_{A/k}\widehat\otimes_A(A/\mathfrak p)
\]
sont fermés pour la topologie $\mathfrak m$-adique ($0_{\mathrm I}$, 7.3.5),
donc la suite (22.7.2.6) est exacte. Tensorisant par $L$, il vient la suite
exacte
\[
  (\mathfrak p/\mathfrak p^2)\otimes_AL\xrightarrow{i}
  \widehat\Omega^1_{A/k}\otimes_AL\longrightarrow
  \widehat\Omega^1_{(A/\mathfrak p)/k}\otimes_{A/\mathfrak p}L
  \longrightarrow0.
\]

Or, l'hypothèse de séparabilité faite sur $L$ entraîne, en vertu de (21.9.5),
que $\widehat\Omega^1_{(A/\mathfrak p)/k}\otimes_{A/\mathfrak p}L$ est de
rang $s$ sur $L$~; comme $\widehat\Omega^1_{A/k}\otimes_AL$ est de rang $r$
sur $L$ (21.9.3), on a $\operatorname{rg}_L(\operatorname{Im}i)=r-s$. Mais comme
$A/\mathfrak p$ est finie sur une sous-algèbre isomorphe à
$k[[T_1,\ldots,T_s]]$, on a $\dim(A/\mathfrak p)=s$ ((17.1.4) et (16.1.5))~;
donc
\[
  r-s=\dim(A)-\dim(A/\mathfrak p)=\dim(A_{\mathfrak p})
\]

